%----------------------------------------------------------------------------
\chapter{The two algorithms}
\label{chap:Algorithms}
%----------------------------------------------------------------------------

\begin{itemize}
    \item This chapter states XPBD and VBD in the form actually implemented in \autoref{chap:Methodology}.
    \item Both target the same variational problem \eqref{eq:G} but reach it through opposite descent directions.
    \item The contrast made explicit here drives the experimental design in \autoref{chap:Experiments}.
\end{itemize}

%----------------------------------------------------------------------------
\section{XPBD: dual-variable position projection}
\label{sec:XPBDAlgo}
%----------------------------------------------------------------------------

\subsection{Derivation}
\label{ssec:XPBDDeriv}
\begin{itemize}
    \item For each elastic spring, associate a constraint $C(\mathbf{x}) = 0$ at zero deformation and a compliance $\alpha = 1/k$.
    \item The augmented Lagrangian of \eqref{eq:G} with a single substep $h$ has stationarity conditions:
\end{itemize}
\begin{align}
    \mathbf{M}(\mathbf{x} - \mathbf{y}) &= h^{2}\,\nabla C(\mathbf{x})^{\!\top}\,\boldsymbol{\lambda}, \\
    C(\mathbf{x}) + \tilde{\alpha}\,\boldsymbol{\lambda} &= \mathbf{0},
\end{align}
\begin{itemize}
    \item $\tilde{\alpha} = \alpha / h^{2}$ is the time-step-scaled compliance.
    \item XPBD solves this coupled system by a Gauss--Seidel sweep over constraints: for each constraint $j$, the multiplier increment is
\end{itemize}
\begin{equation}
    \label{eq:XPBDLambda}
    \Delta \lambda_{j} \;=\; \frac{-C_{j}(\mathbf{x}) - \tilde{\alpha}_{j}\, \lambda_{j}}
                                {\nabla C_{j}^{\!\top}\,\mathbf{M}^{-1}\nabla C_{j}
                                 + \tilde{\alpha}_{j}},
\end{equation}
\begin{itemize}
    \item Corresponding position correction: $\Delta \mathbf{x} = \mathbf{M}^{-1}\nabla C_{j}^{\!\top}\,\Delta \lambda_{j}$.
    \item Velocities are recovered at substep end by $\mathbf{v}_{t+h} = (\mathbf{x}_{t+h} - \mathbf{x}_{t}) / h$.
\end{itemize}

\subsection{Algorithm}
\label{ssec:XPBDPseudo}
\begin{lstlisting}[language=Python,caption={XPBD substep with $n_{\mathrm{iter}}$ Gauss--Seidel iterations.},label=lst:XPBDPseudo]
def xpbd_substep(x, v, h, constraints):
    y = x + h * v + h**2 * M_inv @ f_ext
    x_new = y.copy()
    lambdas = zeros(len(constraints))
    for it in range(n_iter):
        for j, C_j in enumerate(constraints):
            dlam = (-C_j.value(x_new) - alpha_tilde[j] * lambdas[j]) / \
                   (C_j.grad(x_new) @ M_inv @ C_j.grad(x_new).T
                    + alpha_tilde[j])
            x_new += M_inv @ C_j.grad(x_new).T * dlam
            lambdas[j] += dlam
    v_new = (x_new - x) / h
    return x_new, v_new
\end{lstlisting}

\subsection{Characteristics}
\label{ssec:XPBDChar}
\begin{itemize}
    \item \emph{Dual-variable.} An iteration's state is encoded in the multipliers $\boldsymbol{\lambda}$; positions are derived.
    \item \emph{Stiffness-by-compliance.} Material stiffness enters only through $\tilde{\alpha} = 1/(k h^{2})$, so very stiff constraints naturally degenerate to projection ($\tilde\alpha \to 0$).
    \item \emph{First-order gradient information only.} Each constraint contributes its gradient $\nabla C_j$; no Hessian of $E$ is built.
    \item \emph{Substeps preferred over iterations.} Following the small-steps result, I use many short substeps with one inner iteration.
\end{itemize}

%----------------------------------------------------------------------------
\section{VBD: primal block-coordinate descent}
\label{sec:VBDAlgo}
%----------------------------------------------------------------------------

\subsection{Derivation}
\label{ssec:VBDDeriv}
\begin{itemize}
    \item VBD performs block-coordinate descent on $G$ in \eqref{eq:G} with one block per vertex.
    \item Holding all other vertices fixed, the update for vertex $i$ is the Newton step on the local 3-vector problem:
\end{itemize}
\begin{equation}
    \mathbf{x}_{i} \;\leftarrow\; \mathbf{x}_{i} - \mathbf{H}_{i}^{-1}\,\mathbf{g}_{i},
\end{equation}
where
\begin{align}
    \mathbf{g}_{i} &= \frac{m_{i}}{h^{2}}\,(\mathbf{x}_{i} - \mathbf{y}_{i})
                      + \sum_{e \in \mathcal{N}(i)} \nabla_{i} E_{e}(\mathbf{x}), \\
    \mathbf{H}_{i} &= \frac{m_{i}}{h^{2}}\,\mathbf{I}_{3}
                      + \sum_{e \in \mathcal{N}(i)} \nabla_{ii}^{2}\, E_{e}(\mathbf{x}).
\end{align}
\begin{itemize}
    \item For a distance-spring energy $E_{e} = \tfrac{1}{2} k\,(\|\mathbf{x}_{i} - \mathbf{x}_{j}\| - \ell_{0})^{2}$, the per-vertex gradient and Hessian are closed-form and inexpensive.
\end{itemize}

\subsection{Algorithm}
\label{ssec:VBDPseudo}
\begin{lstlisting}[language=Python,caption={VBD step with $n_{\mathrm{iter}}$ block iterations.},label=lst:VBDPseudo]
def vbd_step(x, v, h):
    y = x + h * v + h**2 * M_inv @ f_ext
    x_new = y.copy()  # inertial initialisation
    for it in range(n_iter):
        for i in vertex_order:
            g_i = (m[i] / h**2) * (x_new[i] - y[i])
            H_i = (m[i] / h**2) * I3
            for e in incident_edges(i):
                g_i += spring_grad_i(x_new, e)
                H_i += spring_hess_ii(x_new, e)
            x_new[i] -= solve_3x3(H_i, g_i)
    v_new = (x_new - x) / h
    return x_new, v_new
\end{lstlisting}

\subsection{Characteristics}
\label{ssec:VBDChar}
\begin{itemize}
    \item \emph{Primal.} An iteration's state is positions $\mathbf{x}$; no auxiliary multipliers.
    \item \emph{Second-order locally.} Each block update is a Newton step on the true elastic Hessian restricted to one vertex.
    \item \emph{Iterations over substeps.} VBD is normally used with one long step subdivided by many block iterations; the implicit Hessian handles stiffness without tiny $h$.
    \item \emph{Parallel-friendly.} Vertices sharing no constraint can be updated independently (graph colouring). I disable parallelism for the core comparison; see \autoref{ssec:Parallelisation}.
\end{itemize}

%----------------------------------------------------------------------------
\section{Side-by-side: where the algorithms differ}
\label{sec:Sidebyside}
%----------------------------------------------------------------------------

\begin{table}[!ht]
    \footnotesize
    \centering
    \begin{tabular}{l l l}
        \toprule
                                  & XPBD                          & VBD \\
        \midrule
        Variable updated          & multipliers $\boldsymbol{\lambda}$ & positions $\mathbf{x}$ \\
        Order of information      & gradient $\nabla C$           & gradient + Hessian of $E$ \\
        Block size                & per-constraint                & per-vertex \\
        Stiffness handled by      & compliance $\tilde{\alpha}$   & implicit Hessian \\
        Typical configuration     & many substeps, $1$ iteration  & one step, many iterations \\
        Parallelism strategy      & constraint colouring (harder) & vertex colouring (easier) \\
        \bottomrule
    \end{tabular}
    \caption{Algorithmic contrast between XPBD and VBD.}
    \label{tab:XPBDvsVBD}
\end{table}

\begin{itemize}
    \item The row of \autoref{tab:XPBDvsVBD} that drives the rest of the thesis is the first: XPBD lives in the dual, VBD lives in the primal.
    \item This is the distinction the primal/dual lens (\autoref{sec:PrimalDualLens}) builds on, and the one the experiments return to when interpreting each result.
\end{itemize}
